\documentclass{article}
\usepackage{mathtools} 
\usepackage{fontspec}
\usepackage[UTF8]{ctex}
\usepackage{amsthm}
\usepackage{mdframed}
\usepackage{xcolor}
\usepackage{amssymb}
\usepackage{amsmath}


% 定义新的带灰色背景的说明环境 zremark
\newmdtheoremenv[
  backgroundcolor=gray!10,
  % 边框与背景一致，边框线会消失
  linecolor=gray!10
]{zremark}{说明}

% 通用矩阵命令: \flexmatrix{矩阵名}{元素符号}{行数}{列数}
\newcommand{\flexmatrix}[4]{
  \[
  #1 = \begin{pmatrix}
    #2_{11}     & #2_{12}     & \cdots & #2_{1#4}   \\
    #2_{21}     & #2_{22}     & \cdots & #2_{2#4}   \\
    \vdots      & \vdots      & \ddots & \vdots     \\
    #2_{#31}    & #2_{#32}    & \cdots & #2_{#3#4}
  \end{pmatrix}
  \]
}

% 简化版命令（默认矩阵名为A，元素符号为a）: \quickmatrix{行数}{列数}
\newcommand{\quickmatrix}[2]{\flexmatrix{A}{a}{#1}{#2}}

\begin{document}
\title{注释 7.1}
\author{张志聪}
\maketitle

\begin{zremark}
  推论的证明
\end{zremark}

\textbf{证明：}

通过
\begin{align*}
  \lim\limits_{n \to \infty} a_n = \xi \\
  \lim\limits_{n \to \infty} b_n = \xi
\end{align*}
以证得该推论。

\begin{zremark}
  例3，证明的第一段中$k_\alpha$的存在性。
\end{zremark}
\textbf{证明：}
设$S$的上界为$M$，
阿基米德性保证存在正整数$k$，对任意$s \in S$，有
\begin{align*}
 -M \leq -k \alpha \leq s \leq k \alpha \leq M
\end{align*}

我们考虑$(k - 1) \alpha$，如果存在$\alpha'$使得
\begin{align*}
  \alpha' > (k - 1) \alpha
\end{align*}
则找到$k_\alpha = k$，如果不存在，则$(k - 1)\alpha$是$S$的上界.

继续考虑$(k - 2) \alpha$，如果存在$\alpha'$使得
\begin{align*}
  \alpha' > (k - 2) \alpha
\end{align*}
则找到$k_\alpha = k - 1$，如果不存在，则$(k - 2)\alpha$是$S$的上界.

因为$k$是定值且$S$不是空集，进过有限次上述操作，即可找到$k_\alpha$。

\end{document}